\subsubsection{线性常微分方程组与渐近稳定性}\label{subsec:linear-asymptotic-stability}

现在, 我们在线性常微分方程组的背景下考察一个称为渐近稳定性的概念. 考虑求函数 $\bm{y}:\mathbb{R}\to\mathbb{R}^n$ 的问题, 要求它满足
\begin{equation}\label{eq:linear-nonhomogeneous-system}
\dfrac{\mathrm{d}\bm{y}}{\mathrm{d}t}(t)=\bm{A}(t)\bm{y}(t)+\bm{f}(t),
\end{equation}

\begin{equation}\label{eq:linear-system-initial-data}
\bm{y}(t_0)=\bm{y}_0,
\end{equation}
其中 $t_0\in\mathbb{R}$, $\bm{y}_0\in\mathbb{R}^n$, 向量值函数 $\bm{f}:\mathbb{R}\to\mathbb{R}^n$ 和矩阵值函数 $\bm{A}:\mathbb{R}\to\mathbb{R}^{n\times n}$ 均为给定的.

渐近稳定性描述的是当 $t$ 趋于无穷时齐次方程组的解的行为.

\begin{Definition}
\label{def:linear-system-asymptotic-stability}
称线性齐次方程组
\begin{equation}\label{eq:linear-homogeneous-system}
\bm{y}'=\bm{A}(t)\bm{y}
\end{equation}

\begin{enumerate}
\item 为\emph{渐近稳定的}, 如果 \eqref{eq:linear-homogeneous-system} 的每个解都满足
\begin{equation}\label{eq:asymptotic-stability-limit}
\lim_{t\to\infty}|\bm{y}(t)|=0,
\end{equation}

\item 为\emph{完全不稳定的}, 如果 \eqref{eq:linear-homogeneous-system} 的每个非零解都满足
\begin{equation}\label{eq:complete-instability-limit}
\lim_{t\to\infty}|\bm{y}(t)|=\infty.
\end{equation}
\end{enumerate}
\end{Definition}

下面这个基本结果适用于常系数方程组.

\begin{Theorem}
\label{thm:constant-coefficient-stability}
设 $\bm{A}\in\mathbb{R}^{n\times n}$ 为常矩阵, 其特征值为
\[
\lambda_1,\lambda_2,\ldots,\lambda_n.
\]
那么, 线性齐次常微分方程组
\begin{equation}\label{eq:constant-coefficient-homogeneous-system}
\bm{y}'=\bm{A}\bm{y}
\end{equation}

\begin{enumerate}
\item 渐近稳定, 当且仅当 $\bm{A}$ 的所有特征值的实部都为负; 并且
\item 完全不稳定, 当且仅当 $\bm{A}$ 的所有特征值的实部都为正.
\end{enumerate}
\end{Theorem}

该定理的证明基于矩阵 $\bm{A}$ 的一个对角化过程, 以及与 \eqref{eq:constant-coefficient-homogeneous-system} 对应的初值问题的所有解的下列公式
\begin{equation}\label{eq:matrix-exponential-solution}
\bm{y}(t):=e^{\bm{A}(t-t_0)}\bm{y}_0.
\end{equation}
这里, 矩阵 $e^{\bm{A}t}$ 由一致收敛的幂级数定义为
\begin{equation}\label{eq:matrix-exponential-series}
e^{\bm{A}t}:=\sum_{n=0}^{\infty}\dfrac{\bm{A}^nt^n}{n!}.
\end{equation}


公式 \eqref{eq:matrix-exponential-solution} 是我们将在 Chapter~\ref{chap:semigroup-methods} 中遇到的半群理论中一些公式的先导.

\subsubsection{适定问题}\label{subsec:wellposed-problems}

称一个问题是\emph{适定的 (在 Hadamard 意义下)}, 如果
\begin{enumerate}
\item 存在解,
\item 解是唯一的,
\item 解连续依赖于数据.
\end{enumerate}

如果这些条件不成立, 就称该问题是\emph{不适定的}. 当然, 对于所考察的每个具体问题, 都必须通过选取范数, 使\emph{关于数据的连续性}这一说法的含义更加精确.

在本书中, 我们会将遇到的大多数问题划分为适定问题或不适定问题, 但读者不应认为, 适定问题总比不适定问题"更好"或"更符合物理实际". 正如前面提到的梁的屈曲问题所显示的, 有时适定问题的条件 (在这个例子中是唯一性) 并不符合物理实际. Maxwell 很早就强调了自然界中不适定性的重要性~\cite{Max}:

\begin{quote}
例如, 被冰霜松动而平衡在山坡某个奇点上的岩石, 点燃大片森林的一点火星, 使世界陷入争斗的一句话, 阻止一个人随心所欲的一丝顾忌, 使所有马铃薯枯萎的一点孢子, 使我们成为哲学家或白痴的一点芽球. 凡等级达到某种程度以上的存在, 都有其奇点: 等级越高, 奇点就越多. 在这些点上, 某些影响的物理量值小到有限的存在者无法将其计入, 却可能产生极其重要的结果. 人类努力所产生的一切伟大成果, 都有赖于在这些奇异状态出现时加以利用.
\end{quote}

我们要提醒读者注意, 这段话的发表, 比人们"发现"利用 $\mathbb{R}^3$ 中的三次曲面所能做的种种奇妙事情, 整整早了一个世纪.

\subsubsection{表示公式}\label{subsec:representations}

证明一个问题存在解, 有一种比其他所有方法都更令人满意的方式: 将解显式地写出来. 除了美感上的优点之外, 解的表示公式还有许多实际优点. 利用公式, 可以对解进行计算, 作图, 观察, 估计, 运算和修改. 下面考察一些在偏微分方程研究中经常有用的解的表示公式.

\paragraph*{参数变易法}

\PDEIndexMark{idx:variation-motivation}参数变易公式利用齐次问题 \eqref{eq:linear-homogeneous-system} 的解, 给出非齐次线性常微分方程组 \eqref{eq:linear-nonhomogeneous-system} 的解. 尽管这种表示在实际计算解方面至少有一定用处, 但它主要用于理论分析.

常数变易公式的关键, 是为线性齐次方程组构造一个\emph{基本解矩阵} $\Phi(t,\tau)\in\mathbb{R}^{n\times n}$. 这一解矩阵满足
\begin{equation}\label{eq:fundamental-matrix-dynamics}
\dfrac{\mathrm{d}}{\mathrm{d}t}\Phi(t,\tau)=\bm{A}(t)\Phi(t,\tau),
\end{equation}

\begin{equation}\label{eq:fundamental-matrix-normalization}
\Phi(\tau,\tau)=\bm{I},
\end{equation}
其中 $\bm{I}$ 为 $n\times n$ 单位矩阵. 基本矩阵存在性的证明是标准的, 留作习题.

注意, 齐次方程组的初值问题 \eqref{eq:linear-homogeneous-system}, \eqref{eq:linear-system-initial-data} 的唯一解由下式给出
\begin{equation}\label{eq:fundamental-matrix-homogeneous-solution}
\bm{y}(t):=\Phi(t,t_0)\bm{y}_0.
\end{equation}


利用 Leibniz 公式可知, \emph{参数变易公式}
\begin{equation}\label{eq:variation-of-parameters}
\bm{y}(t):=\Phi(t,t_0)\bm{y}_0+\int_{t_0}^t\Phi(t,s)\bm{f}(s)\,\mathrm{d}s
\end{equation}
给出了非齐次方程组的初值问题 \eqref{eq:linear-nonhomogeneous-system}, \eqref{eq:linear-system-initial-data} 的解.

\paragraph*{Cauchy 积分公式}

Cauchy 积分公式是复变函数理论中最重要的结果. 它利用解析函数在远处各点的值给出该函数的一种表示. 注意, 这种表示很少用于实际计算解析函数的值; 相反, 它被用于推导各种理论结果.

\begin{Theorem}{Cauchy 积分公式}
\label{thm:cauchy-integral-formula}
设 $f$ 在单连通区域 $D\subset\mathbb{C}$ 中解析, 且 $C$ 是 $D$ 中的一条正向简单闭曲线. 那么, 对于 $C$ 内部的任意点 $z_0$, 都有
\begin{equation}\label{eq:cauchy-integral-formula}
f(z_0)=\dfrac{1}{2\pi i}\oint_C\dfrac{f(z)}{z-z_0}\,\mathrm{d}z.
\end{equation}

\end{Theorem}

\subsubsection{估计}\label{subsec:estimation}

说到解的一个\emph{估计}, 我们指的是一种能够表明解的大小或性质的关系. 这些关系大多是涉及解的范数的不等式. 我们区分以下两类估计. \emph{后验估计}依赖于对解的存在性的了解. 这种了解通常通过某种构造或显式表示获得. \emph{先验估计}则是以解的存在为条件的估计; 即形如"如果问题的解存在, 那么它满足 ..."的结果. 这里分别给出这两类估计的一个例子.

\paragraph*{Gronwall 不等式与能量估计}

\PDEIndexMark{idx:gronwall-energy-heading}本节推导常微分方程解的一个先验估计, 它与后面各章将考察的偏微分方程能量估计有关. 唯一性结果 Theorem~\ref{thm:ode-uniqueness} 是这一结果的直接推论. 为了推导所需的估计, 我们需要一个称为 Gronwall 不等式的基本不等式.

\begin{Lemma}{Gronwall 不等式}
\label{lem:gronwall}
设
\[
\begin{aligned}u&:[a,b]\to[0,\infty),\\v&:[a,b]\to\mathbb{R},\end{aligned}
\]
为连续函数, 且 $C$ 为常数. 那么, 如果对 $t\in[a,b]$ 有
\begin{equation}\label{eq:gronwall-integral-hypothesis}
v(t)\leq C+\int_a^tv(s)u(s)\,\mathrm{d}s
\end{equation}
则对 $t\in[a,b]$ 有
\begin{equation}\label{eq:gronwall-exponential-bound}
v(t)\leq C\exp\left(\int_a^tu(s)\,\mathrm{d}s\right)
\end{equation}

\end{Lemma}

这一结果的证明留作习题.

\begin{Lemma}{常微分方程的能量估计}
\label{lem:ode-energy-estimate}
设 $\bm{F}:\mathbb{R}\times\mathbb{R}^n\to\mathbb{R}^n$ 满足 Theorem~\ref{thm:ode-existence} 的假设, 特别地, 它关于第二个变量一致 Lipschitz 连续, Lipschitz 常数为 $\gamma$ (参见 \eqref{eq:uniform-lipschitz-bound}). 设 $\bm{y}_1$ 和 $\bm{y}_2$ 是 \eqref{eq:ode-initial-value-dynamics} 在区间 $[t_0,T]$ 上的解; 即对于 $i=1,2$ 和 $t\in[t_0,T]$, 有
\[
\bm{y}_i'(t)=\bm{F}(t,\bm{y}_i(t))
\]
那么
\begin{equation}\label{eq:ode-energy-estimate}
|\bm{y}_1(t)-\bm{y}_2(t)|^2\leq|\bm{y}_1(t_0)-\bm{y}_2(t_0)|^2e^{2\gamma(t-t_0)}.
\end{equation}

\end{Lemma}

\begin{Proof}
首先利用微分方程, Cauchy-Schwarz 不等式和 Lipschitz 条件, 推导下列不等式.
\[
\begin{aligned}
&|\bm{y}_1(t)-\bm{y}_2(t)|^2\\
&=|\bm{y}_1(t_0)-\bm{y}_2(t_0)|^2+\int_{t_0}^t\dfrac{\mathrm{d}}{\mathrm{d}s}|\bm{y}_1(s)-\bm{y}_2(s)|^2\,\mathrm{d}s\\
&=|\bm{y}_1(t_0)-\bm{y}_2(t_0)|^2\\
&\quad+\int_{t_0}^t2(\bm{y}_1(s)-\bm{y}_2(s))\cdot(\bm{F}(s,\bm{y}_1(s))-\bm{F}(s,\bm{y}_2(s)))\,\mathrm{d}s\\
&\leq|\bm{y}_1(t_0)-\bm{y}_2(t_0)|^2+\int_{t_0}^t2|\bm{y}_1(s)-\bm{y}_2(s)|\,|\bm{F}(s,\bm{y}_1(s))-\bm{F}(s,\bm{y}_2(s))|\,\mathrm{d}s\\
&\leq|\bm{y}_1(t_0)-\bm{y}_2(t_0)|^2+\int_{t_0}^t2\gamma|\bm{y}_1(s)-\bm{y}_2(s)|^2\,\mathrm{d}s.
\end{aligned}
\]
现在, 由 Gronwall 不等式立即得到 \eqref{eq:ode-energy-estimate}.
\end{Proof}

注意, 只需令 $\bm{y}_1(t_0)=\bm{y}_2(t_0)$ 并利用 \eqref{eq:ode-energy-estimate}, 就能得到常微分方程的唯一性结果 (Theorem~\ref{thm:ode-uniqueness}). 还要注意, 这些结果确实是\emph{先验地}得到的: 我们所做的一切都不依赖于解的存在, 而只依赖于解如果\emph{确实}存在就应当满足的方程.

\paragraph*{解析函数的最大值原理}

作为一个\emph{后验}结果的例子, 我们考察下面的定理.

\begin{Theorem}{最大模原理}
\label{thm:maximum-modulus}
设 $D\subset\mathbb{C}$ 为有界区域, 且 $f$ 在 $D$ 上解析, 在 $D$ 的闭包上连续. 那么, $|f|$ 在 $D$ 的边界上取得最大值; 即存在 $z_0\in\partial D$ 使得
\begin{equation}\label{eq:analytic-maximum-modulus}
|f(z_0)|=\sup_{z\in\overline{D}}|f(z)|.
\end{equation}

\end{Theorem}

建议读者利用 Cauchy 积分公式证明这一结果 (参见 Problem~\ref{ex:maximum-modulus-from-cauchy}). 这样的证明基于函数 $f$ 的显式表示, 因而是\emph{后验的}. 但也可以给出这一结果的\emph{先验证明}; Chapter~\ref{chap:maximum-principles} 将专门研究偏微分方程的\emph{先验}最大值原理.

\subsubsection{光滑性}\label{subsec:smoothness}

证明偏微分方程存在解的最重要的现代技巧之一, 是下面这一过程.
\begin{enumerate}
\item 将原来的偏微分方程转化为一种"弱"形式, 这种形式或许可以有很粗糙的解.
\item 证明弱问题存在解.
\item 证明弱方程的解实际上具有比最初预期更多的光滑性.
\item 证明弱问题的一个"光滑"解就是原问题的解.
\end{enumerate}

下文的 Section~\ref{subsec:laplace-equation} 将预先展示这一过程的第一, 第二和第四部分; 而在本节, 我们来考察第三部分, 即证明光滑性的方法的一些先导结果.

\paragraph*{常微分方程解的光滑性}

下面给出一个正则性的"自举"证明的例子: 利用 $\bm{y}\in C^0$ 这一事实证明 $\bm{y}\in C^1$, 再依次类推. 注意, 这一结果可以用来证明 Theorem~\ref{thm:ode-existence} 中关于正则性的部分 (该定理断言存在一个 $C^1$ 解).

\begin{Theorem}
\label{thm:ode-bootstrap-regularity}
如果对某个整数 $m\geq1$, 有 $\bm{F}:\mathbb{R}\times\mathbb{R}^n\to\mathbb{R}^n$ 属于 $C^{m-1}(\mathbb{R}\times\mathbb{R}^n)$, 且 $\bm{y}\in C^0(\mathbb{R})$ 满足积分方程
\begin{equation}\label{eq:ode-bootstrap-integral-equation}
\bm{y}(t)=\bm{y}(t_0)+\int_{t_0}^t\bm{F}(s,\bm{y}(s))\,\mathrm{d}s,
\end{equation}
那么实际上 $\bm{y}\in C^m(\mathbb{R})$.
\end{Theorem}

\begin{Proof}
由于 $\bm{F}(s,\bm{y}(s))$ 连续, 可利用微积分基本定理推出 \eqref{eq:sturm-liouville-kernel-eigenproblem} 的右端连续可微, 因而左端也必须连续可微, 并且
\begin{equation}\label{eq:ode-bootstrap-differential-equation}
\bm{y}'(t)=\bm{F}(t,\bm{y}(t)).
\end{equation}
于是 $\bm{y}\in C^1(\mathbb{R})$. 如果 $\bm{F}$ 属于 $C^1$, 就可以重复这一过程: 注意到 \eqref{eq:ode-bootstrap-differential-equation} 的右端可微 (从而左端也可微), 并且
\[
\bm{y}''(t)=\bm{F}'_{\bm{y}}(t,\bm{y}(t))\cdot\bm{y}'(t)+\bm{F}_t(t,\bm{y}(t)),
\]
所以 $\bm{y}\in C^2(\mathbb{R})$. 只要还能继续对 $\bm{F}$ 求连续导数, 就可以重复上述过程. 因此, 一般而言, $\bm{y}$ 的可微阶数比 $\bm{F}$ 多一阶.
\end{Proof}

\paragraph*{解析函数的光滑性}

对于解析函数, 可以利用 Cauchy 积分公式得到更强的结果.

\begin{Theorem}
\label{thm:analytic-infinite-differentiability}
如果函数 $f:\mathbb{C}\to\mathbb{C}$ 在 $z_0\in\mathbb{C}$ 处解析 (即它在 $z_0$ 的某个邻域内至少有一阶复导数), 那么它具有任意阶复导数. 事实上, 有
\begin{equation}\label{eq:cauchy-higher-derivatives}
f^{(n)}(z_0)=\dfrac{n!}{2\pi i}\oint_C\dfrac{f(z)}{(z-z_0)^{n+1}}\,\mathrm{d}z
\end{equation}
其中 $C$ 为任意一条正向简单闭曲线, 它位于 $f$ 解析的某个单连通区域内, 且 $z_0$ 位于其内部.
\end{Theorem}

在积分号下对 Cauchy 积分公式 \eqref{eq:cauchy-integral-formula} 求导, 就可得到证明. 这是偏微分方程中的一种常用技巧, 读者应当熟悉它 (参见 Problem~\ref{ex:differentiation-under-integral}).

\paragraph*{Problems}

\begin{enumerate}[label=\thesection.\arabic*.,ref=\thesection.\arabic*,series=book-problems]

\item\label{ex:picard-uniform-convergence} 设 $\bm{y}_k$ 为 \eqref{eq:picard-iteration} 定义的序列. 证明
\[
|\bm{y}_{k+1}(t)-\bm{y}_k(t)|\leq\gamma(t-t_0)\max_{\tau\in[t_0,t]}|\bm{y}_k(\tau)-\bm{y}_{k-1}(\tau)|.
\]
利用这一结果证明, 对任意 $T<t_0+1/\gamma$, 该序列在 $t_0\leq t\leq T$ 上一致收敛.

\item\label{ex:sphere-implicit-function} 利用隐函数定理, 确定在什么情况下方程
\[
x^2+y^2+z^2=1
\]
隐式定义函数 $\hat{x}(y,z)$. 给出这一结果的几何解释.

\item\label{ex:implicit-function-regularity} 证明, 如果 Theorem~\ref{thm:implicit-function} 中的 $\bm{F}$ 为 $C^k$ 的, 那么 $\hat{\bm{y}}$ 也为 $C^k$ 的. Hint: 首先考虑差商, 证明 $\hat{\bm{y}}$ 为 $C^1$ 的.

\item\label{ex:nonunique-minimizers} 证明, 在满足 $u(0)=u(1)=0$ 的所有分片 $C^1$ 函数所构成的集合上, 泛函
\[
\mathcal{E}(u):=\int_0^1(1-(u'(t))^2)^2\,\mathrm{d}t
\]
有无穷多个极小元.

\item\label{ex:nonattained-minimum} 证明, 泛函
\[
\mathcal{E}(u):=\int_0^1u(t)^2+(1-(u'(t))^2)^2\,\mathrm{d}t
\]
不存在满足 $u(0)=u(1)=0$ 的分片 $C^1$ 极小元. Hint: 利用上一题的解的一个序列, 证明极小元 $\tilde{u}$ 必须满足 $\mathcal{E}(\tilde{u})=0$. Remark: 具有这类特点的极小化问题会出现在相变建模中.

\item\label{ex:nonuniform-initial-data-continuity} 举例说明, 当 $t\to\infty$ 时, Theorem~\ref{thm:initial-data-continuity} 中的 $\delta$ 不能选得与 $t$ 无关.

\item\label{ex:distinct-eigenvalue-stability} 当 $\bm{A}$ 的特征值互异时, 证明 Theorem~\ref{thm:constant-coefficient-stability}.

\item\label{ex:fundamental-matrix-existence} 证明 \eqref{eq:fundamental-matrix-dynamics}, \eqref{eq:fundamental-matrix-normalization} 的解的存在唯一性.

\item\label{ex:gronwall-proof} 证明 Gronwall 不等式.

\item\label{ex:maximum-modulus-from-cauchy} 利用 Cauchy 积分公式证明 Theorem~\ref{thm:maximum-modulus}.

\item\label{ex:differentiation-under-integral} 设 $g:\mathbb{R}^2\to\mathbb{R}$ 为 $C^1$ 的. 定义
\[
f(x):=\int_a^bg(x,y)\,\mathrm{d}y.
\]
利用差商证明, 可以"在积分号下"对 $f$ 求导.

\end{enumerate}

\subsection{初等偏微分方程}\label{sec:elementary-pdes}

上一节讨论了本书其余部分将始终考察的几类基本数学问题, 并考察了偏微分方程的两个分支, 即常微分方程和复变函数理论, 是如何回答这些问题的. 现在来看, 初等偏微分方程课程通常如何处理这些问题. 为此, 我们考察三个基本偏微分方程 (Laplace 方程, 热方程和波动方程). 尽管有时也采用分析方法研究它们的性质, 但我们的基本技巧是显式地计算解. 此时我们并不十分关注严密性或一般性, 而着重预示后面的内容; 这里介绍的所有方法和观察, 都将在后面加以推广.
